\section{Introducción}
El diseño y análisis de redes eléctricas frecuentemente requiere la transición de modelos teóricos o prototipos ideales hacia implementaciones físicas con componentes comerciales. Los prototipos normalizados, aunque útiles para simplificar el cálculo matemático inicial al emplear valores unitarios, no siempre representan las condiciones de operación reales en términos de magnitud de componentes o frecuencias de trabajo.

En esta práctica profundizaremos en las técnicas de escalamiento de circuitos, las cuales permiten modificar los parámetros de una red sin alterar la naturaleza de su comportamiento. A través del uso de herramientas de simulación y análisis matemático, se verificará cómo el escalamiento de magnitud y frecuencia facilita la adaptación de un diseño a requerimientos específicos, manteniendo la integridad de las funciones de transferencia y la respuesta del sistema.

